% Generated draft: requires mathematical and source review.
\documentclass[11pt]{article}
\usepackage[margin=25mm]{geometry}
\usepackage{fontspec}
\setmainfont{Latin Modern Roman}
\usepackage{amsmath,amssymb,mathtools,unicode-math}
\setmathfont{Latin Modern Math}
\usepackage{enumitem,needspace,xurl,hyperref}
\setlength{\parindent}{0pt}
\setlength{\parskip}{5pt}
\setlength{\emergencystretch}{3em}
\newcommand{\N}{\mathbb{N}}
\newcommand{\Z}{\mathbb{Z}}
\newcommand{\Q}{\mathbb{Q}}
\newcommand{\R}{\mathbb{R}}
\newcommand{\C}{\mathbb{C}}
\newcommand{\F}{\mathbb{F}}
\newcommand{\A}{\mathbb{A}}
\newcommand{\PP}{\mathbb{P}}
\newcommand{\E}{\mathbb{E}}
\newcommand{\eps}{\varepsilon}
\newcommand{\dd}{\,\mathrm{d}}
\begin{document}
\section*{UnsolvedMath ID: 7200015}
\textbf{Source problem code:} catalogue entry
\subsection{Definitions and mathematical statement}
For the usual affine formulation let \(A\) be a finitely generated commutative integral \(\C\)-algebra and \(X=\operatorname{Spec}A\). Write \(m_a\) for multiplication by \(a\). Define the ring of \(\C\)-linear differential operators inductively by
\[
D^0_{A/\C}=\{m_a:a\in A\},\qquad
D^{r+1}_{A/\C}=\{\theta\in\operatorname{End}_{\C}(A):
 [\theta,m_a]\in D^r_{A/\C}\text{ for all }a\in A\},
\]
\[
D_{A/\C}=\bigcup_{r\ge0}D^r_{A/\C}.
\]
A derivation \(\delta\in\operatorname{Der}_{\C}(A)\) satisfies \(\delta(ab)=a\delta(b)+b\delta(a)\). Nakai's conjecture in this setting is
\[
D_{A/\C}=\langle m_a\ (a\in A),\ \operatorname{Der}_{\C}(A)\rangle
\quad\Longrightarrow\quad X\text{ is smooth over }\C,
\]
where the brackets mean the generated associative \(\C\)-algebra. Smoothness is equivalently regularity of every local ring of \(A\).

Order-zero multiplication operators must be included in the conventional generation condition; derivations alone annihilate \(1\), so literal generation solely by derivations is not the usual statement. For a nonaffine variety the standard local version uses the sheaf of differential operators generated by \(\mathcal O_X\) and the derivation sheaf. Generation of only its ring of global sections is a different hypothesis. The extraction does not specify affine versus sheaf-local scope or integral versus general reduced varieties, so those extensions remain for source review.

The supplied research summary reports the general conjecture open and results for toric varieties, certain quotients, isolated homogeneous hypersurface singularities, and specified low-modality hypersurfaces. It cites two recent preprints from abstract-level checks, not a general resolution. These reported special cases are not extended to arbitrary homogeneous or arbitrary hypersurface varieties.
\subsection{Short English statement}
If every differential operator on a complex affine variety is built from functions and derivations, must the variety be smooth? The usual order-zero generators and the local versus global formulation must be made explicit.
\subsection{Sources}
\begin{itemize}
\item Local UnsolvedMath registry, permanent ID \(7200015\), exact source code \texttt{catalogue entry}, ``Nakai conjecture.''
\item Wikipedia, \emph{List of unsolved problems in mathematics}, pinned revision 1366636767, Geometry bullet 15; \url{https://en.wikipedia.org/w/index.php?title=List_of_unsolved_problems_in_mathematics&oldid=1366636767#Geometry}. MediaWiki-source extraction; recorded access 2026-07-29.
\item Supplied research summary: ``The Nakai Conjecture for isolated hypersurface singularities of modality \(\le2\),'' arXiv:2502.04672 (2025); ``Nakai conjectures for isolated homogeneous hypersurface singularities,'' arXiv:2604.24508 (2026). Authors and full publication details are not supplied; the summary describes abstract-level checks.
\item Other reported toric/quotient and exceptional-hypersurface special cases have no complete citations in this record. A primary-source general-conjecture citation is not supplied.
\end{itemize}
\end{document}
